\part[Кратные интегралы и теория поля]{Кратные интегралы и теория поля}

\chapter{Формула Грина. Потенциальные векторные поля на плоскости.}
\section{Формула Грина} 

\subsection{Каноническая область и ориентация границы}

\begin{defn}
Пусть $G$ "--- ограниченное открытое множество на плоскости, граница которого является непрерывной кривой (см. рис. 19.17). Граница $\partial G$ считается \textit{положительно ориентированной} (обозначается $\partial G^+$), если при обходе границы множество $G$ остаётся слева, т.е. положительная ориентация соответствует направлению обхода против часовой стрелки. Если граница открытого множества состоит из конечного числа непрерывных кривых, в частности, имеются <<дырки>> (см. рис. 19.18), то положительная ориентация границы <<дырки>> соответствует её обходу по часовой стрелке. Направление обхода $\partial G$, противоположное положительной ориентации, называется \textit{отрицательной ориентацией} ($\partial G^-$).
\end{defn}

\begin{defn}
Плоская область $G$ называется \textit{канонической\rindex{область!каноническая}}, если она одновременно задаётся как неравенствами
$$
a < x < b, \quad \psi(x) < y < \varphi(x),
$$
где функции $\psi$ и $\varphi$ непрерывны на $[a;b]$, так и неравенствами
$$
c < y < d, \quad \xi(y) < x < \eta(y),
$$
где функции $\xi$ и $\eta$ непрерывны на $[c;d]$.

Замыкание канонической области является одновременно простым множеством как относительно оси $y$, так и относительно оси $x$.
\end{defn}

\begin{defn}
Пусть $G$ "--- ограниченная область в $\bbR^n$. Функция $f$ называется $k$ \textit{раз непрерывно дифференцируемой в} $\overline{G}$ (обозначается $f \in C^k(\overline{G})$), если она непрерывна в $\overline{G}$, $k$ раз непрерывно дифференцируема в $G$, и все её частные производные порядков $1, 2, \ldots, k$ продолжаются на $\overline{G}$ до функций, непрерывных в $\overline{G}$. При $k = 1$ такая функция называется \textit{непрерывно дифференцируемой в} $\overline{G}$.
\end{defn}

\subsection{Формула Грина}

Формула Грина связывает значение криволинейного интеграла второго рода по границе плоской области со значением некоторого двойного интеграла по всей области.

\begin{thm}[формула Грина]
Пусть $G$ "--- каноническая область на плоскости, функции $P(x,y)$ и $Q(x,y)$ непрерывно дифференцируемы в $\overline{G}$. Тогда
$$
\iint\limits_{G} \left( \pd{Q}{x} - \pd{P}{y} \right) dx\,dy = \int\limits_{\partial G^+} \bigl( P(x,y)\,dx + Q(x,y)\,dy \bigr).
$$
\end{thm}

\begin{notion}
Доказательство теоремы проходит для более общего случая, когда $P$ и $Q$ непрерывны в $\overline{G}$, $\pd{P}{y}$ и $\pd{Q}{x}$ непрерывны и ограничены в $G$, причём существуют конечные пределы
\begin{equation} \label{Grin_lim}
\begin{aligned}
&\lim_{y \to \varphi(x)-0} \pd{P}{y},\quad \lim_{y \to \psi(x)+0} \pd{P}{y}, \quad x \in (a;b), \\
&\lim_{x \to \eta(y)-0} \pd{Q}{x},\quad \lim_{x \to \xi(y)+0} \pd{Q}{x}, \quad y \in (c;d)
\end{aligned}
\end{equation}
(обозначения как в определении 19.8).
\end{notion}

\begin{proof}
Каноническая область на плоскости имеет вид, изображённый на рис. 19.19. Кривая $D_2 A_1 A_2 B_1$ "--- график функции $y = \varphi(x)$, кривая $D_1 C_2 C_1 B_2$ "--- график функции $y = \psi(x)$. Так как $\psi(x) < \varphi(x)$ при $x \in (a;b)$, а функции $\psi$ и $\varphi$ непрерывны на $[a;b]$, то $\psi(a) \leqslant \varphi(a)$, $\psi(b) \leqslant \varphi(b)$. Точка $D_2$ расположена выше $D_1$ или совпадает с ней, точка $B_1$ расположена выше $B_2$ или совпадает с ней. Если точка $D_2$ выше $D_1$, то отрезок $D_1 D_2$ целиком принадлежит границе области $G$, если точка $B_1$ выше $B_2$, то отрезок $B_1 B_2$ целиком принадлежит границе $G$. Таким образом, $\partial G$ состоит из графиков непрерывных функций $D_2 A_1 A_2 B_1$ и $D_1 C_2 C_1 B_2$ и отрезков $D_1 D_2$ и $B_1 B_2$, т.е. $\partial G$ имеет меру нуль в $\bbR^2$. Следовательно, область $G$ измерима в $\bbR^2$. Пока мы учли только то, что $\overline{G}$ "--- простое множество относительно оси $y$. Если учесть то же самое относительно оси $x$, то мы увидим, что границе $G$ принадлежат отрезки $A_1 A_2$ и $C_1 C_2$.

Рассмотрим $\iint\limits_G \pd{P}{y}\,dx\,dy$. Он существует по теореме 19.5 (в более общем случае по теореме 19.9). Так как $\mu \partial G = 0$, то, применяя сведение двойного интеграла к повторному, имеем
$$
\iint\limits_G \pd{P}{y}\,dx\,dy = \iint\limits_{\overline{G}} \pd{P}{y}\,dx\,dy = \int_a^b dx \int_{\psi(x)}^{\varphi(x)} \pd{P(x,y)}{y}\,dy.
$$

При фиксированном $x \in (a;b)$ функция $P(x,y)$ непрерывна по $y$ на отрезке $[\psi(x); \varphi(x)]$, непрерывно дифференцируема на интервале $(\psi(x); \varphi(x))$, и $\pd{P}{y}$ имеет конечные пределы в концах интервала. С учётом теоремы 4.16 эта функция непрерывно дифференцируема на $[\psi(x); \varphi(x)]$, и по формуле Ньютона"--~Лейбница
$$
\int_{\psi(x)}^{\varphi(x)} \pd{P(x,y)}{y}\,dy = P(x,\varphi(x)) - P(x,\psi(x)).
$$

Тогда
\begin{multline*}
\iint\limits_G \pd{P}{y}\,dx\,dy = \int_a^b P(x,\varphi(x))\,dx - \int_a^b P(x,\psi(x))\,dx = \\
= \int\limits_{\overline{D_2 A_1 A_2 B_1}} P(x,y)\,dx - \int\limits_{\overline{D_1 C_2 C_1 B_2}} P(x,y)\,dx = \\
= -\left( \int\limits_{\overline{B_1 A_2 A_1 D_2}} P\,dx + \int\limits_{\overline{D_1 C_2 C_1 B_2}} P\,dx \right)
\end{multline*}
(см. определение 19.6). Так как
$$
\int\limits_{\overline{D_2 D_1}} P\,dx = \int\limits_{\overline{B_2 B_1}} P\,dx = 0
$$
(по лемме 19.8), то в сумме имеем
$$
\iint\limits_G \pd{P}{y}\,dx\,dy = -\int\limits_{\partial G^+} P\,dx.
$$

Аналогично, $\iint\limits_G \pd{Q}{x}\,dx\,dy = \int\limits_{\partial G^+} Q\,dy$. Смена знака обусловлена тем, что <<при замене $x$ на $y$ меняется ориентация плоскости>>; аккуратно это можно проследить, проводя соответствующие выкладки. Сложив два последних равенства, получим:
$$
\iint\limits_G \left( \pd{Q}{x} - \pd{P}{y} \right) dx\,dy = \int\limits_{\partial G^+} \bigl( P\,dx + Q\,dy \bigr).
$$
\end{proof}

Имеет место следующее обобщение этой теоремы.

\begin{thm}[19.20$'$]
Пусть $G$ "--- ограниченное открытое множество на плоскости, граница $\partial G$ которого состоит из конечного числа кусочно-гладких кривых. При этом $G$ разбивается на конечное число канонических областей $G_1, \ldots, G_N$. Это значит, что $\overline{G_1}, \overline{G_2}, \ldots, \overline{G_N}$ образуют разбиение $\overline{G}$; разные $G_i$ не имеют общих точек, а $\overline{G_i}$ могут иметь общие точки разве что на границах. Тогда
$$
\iint\limits_G \left( \pd{Q}{x} - \pd{P}{y} \right) dx\,dy = \int\limits_{\partial G^+} \bigl( P\,dx + Q\,dy \bigr),
$$
если функции $P(x,y)$ и $Q(x,y)$ непрерывно дифференцируемы в $\overline{G}$ (по поводу ориентации $\partial G^+$ см. определение 19.7).
\end{thm}


\section{Потенциальные векторные поля}

Если на множестве $G \subset \bbR^n$ задана векторная функция $\vec{a}(M)$, то говорят, что на $G$ задано \textit{векторное поле\rindex{векторное поле}} $\vec{a}(M)$.

\begin{defn}
Непрерывно дифференцируемая скалярная функция $u$ называется \textit{потенциалом\rindex{потенциал векторного поля}} непрерывного векторного поля $\vec{a} = (P, Q, R)^T$ в области $G \subset \bbR^3$, если в этой области $\vec{a} = \grad u$ (т.е. $P = \pd{u}{x}$, $Q = \pd{u}{y}$, $R = \pd{u}{z}$). Аналогичное определение даётся для области $G \subset \bbR^2$; там $\vec{a} = (P, Q)^T$, $P = \pd{u}{x}$, $Q = \pd{u}{y}$.
\end{defn}

\begin{defn}
Векторное поле $\vec{a}$ называется \textit{потенциальным\rindex{векторное поле!потенциальное}} в области $G \subset \bbR^3$ (или $\bbR^2$), если оно имеет потенциал в этой области.
\end{defn}

\begin{defn}
Пусть ориентированная кусочно-гладкая кривая в $\bbR^2$ или $\bbR^3$ задаётся уравнением $\vec{r} = \vec{r}(t)$, $t \in [a;b]$, причём точка, соответствующая значению параметра $t = a$ является начальной, а соответствующая значению $t = b$ "--- конечной. Кривая называется \textit{замкнутой\rindex{кривая!замкнутая}}, если начальная и конечная точки её совпадают (сравнить с определением 6.13 неориентированной простой замкнутой кривой).
\end{defn}

\begin{thm}
Пусть вектор-функция $\vec{a}$ непрерывна в области $G \subset \bbR^3$ (или $\bbR^2$). Тогда следующие $3$ утверждения равносильны:
\begin{enumerate}
\item[$1^\circ$.] $\int\limits_\Gamma (\vec{a}, d\vec{r}) = 0$ по любой замкнутой кусочно-гладкой кривой $\Gamma \subset G$.
\item[$2^\circ$.] $\int\limits_\gamma (\vec{a}, d\vec{r})$ по кусочно-гладкой кривой $\gamma \subset G$ зависит только от начальной и конечной точек кривой, но не зависит от самой кривой.
\item[$3^\circ$.] Векторное поле $\vec{a}$ имеет потенциал.
\end{enumerate}
При выполнении каждого из этих условий $\int\limits_\gamma (\vec{a}, d\vec{r}) = u(B) - u(A)$, где $A$ "--- начальная, $B$ "--- конечная точка кривой, $u$ "--- потенциал.
\end{thm}

\begin{proof}
Рассмотрим случай $\bbR^3$ (случай $\bbR^2$ разбирается аналогично).

$1^\circ \Rightarrow 2^\circ$ (см. рис. 21.3). Рассмотрим две кусочно-гладкие кривые $\gamma_1$ и $\gamma_2$ с общим началом $A$ и концом $B$. Они образуют замкнутую ориентированную кусочно-гладкую кривую $\Gamma$ (сначала из точки $A$ движемся к точке $B$ по кривой $\gamma_1$, затем обратно из $B$ в $A$ по кривой $\gamma_2$). Тогда
$$
0 = \int\limits_\Gamma (\vec{a}, d\vec{r}) = \int\limits_{A\gamma_1 B} + \int\limits_{B\gamma_2 A} = \int\limits_{A\gamma_1 B} - \int\limits_{A\gamma_2 B}, \quad \text{т.е.} \quad \int\limits_{A\gamma_1 B} = \int\limits_{A\gamma_2 B}
$$
(интегралы по кривым $\gamma_1$ и $\gamma_2$ с общими началом $A$ и концом $B$ совпадают). Отметим, что кривая $\Gamma$ может иметь точки самопересечения, т.е. не обязана являться простой замкнутой.

$2^\circ \Rightarrow 3^\circ$. Пусть $A = (x_0, y_0, z_0) \in G$ "--- фиксированная точка, $B = (x, y, z) \in G$ "--- произвольная точка. Так как $\int\limits_\gamma (\vec{a}, d\vec{r})$ зависит только от начальной и конечной точек кривой, то выражение
$$
u(x, y, z) = \int\limits_A^B (\vec{a}, d\vec{r})
$$
определяет однозначную функцию точки $B$, не зависящую от кусочно-гладкой кривой $\gamma$ с началом в $A$ и концом в $B$. Докажем, что $u$ "--- потенциал векторного поля $\vec{a}$.

Так как $G$ "--- область, то точка $B$ входит в $G$ вместе с некоторой окрестностью, и для достаточно малых $\tau$ точка $C(x + \tau, y, z)$ и весь отрезок $BC$ целиком принадлежат $G$ (см. рис. 21.4). Тогда
\begin{multline*}
u(x + \tau, y, z) - u(x, y, z) = \int\limits_A^C (\vec{a}, d\vec{r}) - \int\limits_A^B (\vec{a}, d\vec{r}) = \int\limits_B^C (\vec{a}, d\vec{r}) = \\
= \int\limits_B^C P(\xi, \eta, \zeta)\,d\xi + Q(\xi, \eta, \zeta)\,d\eta + R(\xi, \eta, \zeta)\,d\zeta.
\end{multline*}
Этот интеграл не зависит от кусочно-гладкой кривой, соединяющей точки $B$ и $C$, его можно брать по отрезку $BC$. Параметризуем отрезок $BC$: $\xi = t$, $\eta = y$, $\zeta = z$, $t \in [x; x + \tau]$ ($x, y, z$ "--- фиксированные числа, координаты точки $B$; $t$ "--- параметр). Если $\tau < 0$, то отрезок $BC$ всё равно пробегается от $B$ к $C$; в этом случае ориентация отрезка соответствует убыванию $t$. Тогда
$$
u(x + \tau, y, z) - u(x, y, z) = \int\limits_x^{x + \tau} P(t, y, z)\,dt = P(\tau_0, y, z) \cdot \tau,
$$
где $\tau_0 = \tau_0(\tau) \in [x; x + \tau]$; здесь применена теорема о среднем для определённого интеграла по отрезку.

Так как $\lim\limits_{\tau \to 0} \tau_0(\tau) = x$, а функция $P$ непрерывна по первому аргументу, то $\lim\limits_{\tau \to 0} P(\tau_0, y, z) = P(x, y, z)$ и существует
$$
\pd{u}{x}(x, y, z) = \lim\limits_{\tau \to 0} \frac{u(x + \tau, y, z) - u(x, y, z)}{\tau} = \lim\limits_{\tau \to 0} P(\tau_0, y, z) = P(x, y, z).
$$
Аналогично, $\pd{u}{y} = Q$, $\pd{u}{z} = R$. Значит, функция $u$ непрерывно дифференцируема в области $G$ и является потенциалом поля $\vec{a}$.

$3^\circ \Rightarrow 1^\circ$. Пусть $P = \pd{u}{x}$, $Q = \pd{u}{y}$, $R = \pd{u}{z}$ в области $G$; $\gamma$ "--- произвольная кусочно-гладкая кривая в $G$, параметризуемая параметром $t$: $x = x(t)$, $y = y(t)$, $z = z(t)$, $t \in [\alpha; \beta]$; $A = (x(\alpha), y(\alpha), z(\alpha))$ "--- начало, $B = (x(\beta), y(\beta), z(\beta))$ "--- конец кривой. Тогда
\begin{multline*}
\int\limits_\gamma (P\,dx + Q\,dy + R\,dz) = \int\limits_\alpha^\beta \left( \pd{u}{x} x'(t) + \pd{u}{y} y'(t) + \pd{u}{z} z'(t) \right) dt = \\
= \int\limits_\alpha^\beta \frac{d}{dt}\bigl( u(x(t), y(t), z(t)) \bigr)\,dt = \\
= u(x(\beta), y(\beta), z(\beta)) - u(x(\alpha), y(\alpha), z(\alpha)) = u(B) - u(A).
\end{multline*}
Для замкнутой кусочно-гладкой кривой $\Gamma \subset G$ начало и конец совпадают, т.е. $A = B$ и $\int\limits_\Gamma (\vec{a}, d\vec{r}) = 0$.

Попутно доказано, что при наличии потенциала для кривой $\gamma$ с началом $A$ и концом $B$ интеграл равен $u(B) - u(A)$. Равенство $\int\limits_A^B (\vec{a}, d\vec{r}) = u(B) - u(A)$ является аналогом формулы Ньютона"--~Лейбница для потенциальных полей.
\end{proof}

\begin{notion}
Потенциал непрерывного векторного поля определяется однозначно с точностью до прибавления произвольной постоянной $C$.
\end{notion}

\begin{proof}
В самом деле, пусть непрерывная вектор-функция $\vec{a}$ имеет два потенциала $u$ и $v$ в области $G$. Тогда если $A$ "--- фиксированная, а $B$ "--- произвольная точка области $G$, то
$$
\int\limits_A^B (\vec{a}, d\vec{r}) = u(B) - u(A) = v(B) - v(A),
$$
т.е. $u(B) = v(B) + u(A) - v(A)$. Так как $u(A) - v(A) = C$ (постоянная величина), то $u(B) = v(B) + C$ для любой точки $B \in G$.
\end{proof}